\documentclass{memoir}
\usepackage[utf8]{inputenc}
\usepackage[T1]{fontenc}
\usepackage{geometry}
\usepackage{graphicx}
\usepackage{float}
\usepackage{hyperref}
\usepackage[english,indonesian]{babel}

\title{Optimasi Pemrosesan Template LaTeX Menggunakan Pendekatan Berbasis AI}
\author{Academic Support Agent}
\date{\today}

\begin{document}
\maketitle

\begin{abstract}

\end{abstract}

\chapter{PENDAHULUAN}


\chapter{TINJAUAN PUSTAKA}


% Original content (commented out):
% % *************** LEd - dokument oparty na szablonie pracy magisterskiej ***************

\documentclass[12pt, a4paper, oneside]{memoir}

\input{style.tex}

\begin{document}
%	\sffamily
% *************** Front matter ***************
\frontmatter
\input{front.tex}



% *************** Main matter ***************
\mainmatter

\addtocounter{page}{1}
\input{chap1//chap1.tex}
\input{chap2//chap2.tex}
\input{chap3//chap3.tex}
\input{chap4//chap4.tex}

% *************** Bibliography ***************
\bibliographystyle{plain}
\bibliography{MSc}

% *************** Back matter ***************
\backmatter
\input{back.tex}

% *************** Appendixes ***************
\appendix
\input{app1//app1.tex}



\end{document}


\end{document}